\section{Experimental setup}
\label{sec:exp_setup}

\textbf{Candidates.} We analyze (1) the Last Input Token (LIT) from all layers and (2) LIT from the final layer. Hidden states in the last position of the prompt is what the language-model head projects to logits, encoding the model's immediate pre-generation context prior to emitting $y$. Prior work confirms that linear probes on these pre-generation activations recover non-trivial information about the upcoming generation~\citep{lugoloobi2026llms}. We evaluate the two variants as reasoning-relevant computation is often concentrated in middle rather than final layers~\citep{cai2026tmlr}. We additionally evaluate (3) soft tokens with no noise (ST)~\citep{zhang2025softthinkingunlockingreasoning} and with Gumbel noise (STN)~\citep{wu2026llms}, and (4) latent thinking (LT)~\citep{zou2025latentcollaborationmultiagentsystems} (see~\cref{appendix:sub:soft_latent_formula}). For the soft tokens and latent thinking methods, we test varying thinking steps of 1, 16, 32, 64, and 128. We also evaluate exact and pooled output embeddings and the input prompt embedding.

% Result-table column groups, wrapped next to the Result-table
% organisation paragraph. "Group" used as neutral header that covers
% upper-anchor (Output Embedding), under-audit (Candidates), and
% reference (Baselines) rows alike.
\begin{wraptable}{r}{0.38\linewidth}
  \vspace{-0.6em}
  \caption{Result-table column groups.}
  \label{tab:candidate-blocks}
  \centering
  \footnotesize
  \setlength{\tabcolsep}{4pt}
  \renewcommand{\arraystretch}{1.15}
  \begin{tabular}{@{}l l@{}}
    \toprule
    \textbf{Group} & \textbf{Variants} \\
    \midrule
    Output Emb. & Exact, Pooled \\
    Candidates       & Hidden states, Think. methods \\
    Baselines        & IE, RV \\
    \bottomrule
  \end{tabular}
  \vspace{-0.6em}
\end{wraptable}

\textbf{Table layout.} Columns follow \cref{tab:candidate-blocks}. The Output Embedding (OE) block holds two upper-bound references derived from $Y$, the Exact variant carrying direct semantic knowledge of $Y$ and the Pooled variant averaging embeddings of possible generations. Neither variant is a reference for Minimality or Stability, where output-based encodings carry their own penalties. The Input Embedding (IE) is the prompt embedding, so a candidate failing to outperform IE adds no information beyond the prompt. The Random Vector (RV) is an information-free reference point.

% Source LLMs at-a-glance, wrapped next to the Source LLMs paragraph.
% Model names abbreviated per the convention used by the appendix
% headline tables (DS-R1-Qwen 32B, Skywork-OR1 32B, GPT-OSS 20B etc.)
\begin{wraptable}{r}{0.44\linewidth}
  \vspace{-0.6em}
  \caption{Source LLMs covered by the audit.}
  \label{tab:source-llms}
  \centering
  \footnotesize
  \setlength{\tabcolsep}{4pt}
  \renewcommand{\arraystretch}{1.1}
  \begin{tabular}{@{}l l l@{}}
    \toprule
    \textbf{Source LLM} & \textbf{Family} & \textbf{Paradigm} \\
    \midrule
    Llama-3.1 8B    & Dense      & Instruct          \\
    Llama-3.3 70B   & Dense      & Instruct          \\
    DS-R1-Qwen 32B  & Dense      & Reasoning-distill \\
    Skywork-OR1 32B & Dense      & Native RL         \\
    GPT-OSS 20B     & Sparse MoE & Adjust effort  \\
    \bottomrule
  \end{tabular}
  \vspace{-0.6em}
\end{wraptable}

\textbf{Models.} We evaluate on the 23 tasks of BBEH \citep{kazemi2025bigbenchextrahard} using the original benchmark prompt. We use five open-weight LLMs chosen to cover a range of sizes, training procedures, and architectures (\cref{tab:source-llms}), specifically Llama-3.1-8B-Instruct and Llama-3.3-70B-Instruct~\citep{grattafiori2024llama3herdmodels}, DeepSeek-R1-Distill-Qwen-32B~\citep{deepseekai2025deepseekr1incentivizingreasoningcapability}, Skywork-OR1-32B~\citep{he2025skywork, skywork-or1-2025}, and GPT-OSS-20B~\citep{openai2025gptoss120bgptoss20bmodel}. The selection covers dense, sparse-MoE, reasoning-distilled, and RL-trained paradigms.

\textbf{Generation.} We assume that $\mathbf{T}$ encodes the sufficient statistics of $P(Y \mid x)$, so the outputs the model assigns high probability to are those consistent with what $\mathbf{T}$ captures about the input. For each prompt, beam search approximates this high-probability region of $P(Y \mid x)$, and the eight returned sequences of up to 8192 tokens form an empirical representative slice on which each axiom is evaluated. Beyond maximizing output probability, beam search guarantees distinct candidate outputs, reduces sampling variance, and exposes an empirical distribution over reasoning paths~\citep{fadeeva2026dont}.

\textbf{Probes.} We utilize a frozen LlaMA-3.2-1B~\citep{grattafiori2024llama3herdmodels} backbone with a trainable projection that maps thought representations into its token-embedding space, and the discriminator adds a trained classification head. The frozen backbone serves as a shared decoding surface whose learned representations remain compatible with those of independently trained LLMs~\citep{salhan2026do}, while the trainable projection learns features specific to each candidate and source model pair. For the Causality measure specifically, the projection is trained on $\mathcal{M}_\theta$'s own output sequences, thus the KL divergence reflects functional substitution rather than generic transferability. Utilizing a shared backbone ensures that the computational cost of evaluating the metric remains constant, independent of the LLM's size. Semantic similarity between outputs for equivalence classes of DCS are computed with Embed-Nemotron-8B~\citep{babakhin2025llamaembednemotron8buniversaltextembedding}, the leading text embedding model on MTEB~\citep{muennighoff2023mtebmassivetextembedding} at the time of writing. Training parameters and auxiliary details are in \cref{appendix:sub:train_details,appendix:sub:bootstrap,appendix:output_length,appendix:bbeh_accuracy}.

